Pool allocator (пулловый аллокатор) "--- это аллокатор, который выделяет память блоками фиксированного размера (пулами). Его идея заключается в том, что некоторый участок памяти, выделенный аллокатору, разделяется на более мелкие участки одинакового размера. Это значительно ускоряет выделение и освобождение памяти для объектов одного размера, минимизирует фрагментацию памяти и локализует данные в памяти для лучшего попадания в кэш. К минусам пуллового аллокатора можно отнести его неуниверсальность: если возникает необходимость в объектах нескольких разных размеров, требуется создавать отдельные аллокаторы под каждый размер.

При запросе выделения памяти пулловый аллокатор возвращает указатель на один из свободных участков памяти, а при освобождении памяти он запоминает, что теперь этот блок можно заново использовать. Для хранения информации о состоянии блока (занят он или свободен) можно использовать булевый массив или односвязный список. В последнем случае при выделении памяти выполняются следующие шаги \cite{allocators}.

\begin{enumerate}
	\item проверка, остались ли звенья в списке свободных блоков (если их нет, то память в аллокаторе закончилась);
	\item передача адреса первого/последнего звена списка пользователю;
	\item удаление этого звена из списка.
\end{enumerate}

При освобождении блока памяти аллокатор добавляет адрес памяти в начало или конец односвязного списка. (Используются именно начало/конец, т.к. вставка туда производится за константное время, в отличие от произвольного места, куда вставка выполняется за $O(n)$). Однако может возникнуть ситуация, когда пользователь ошибся адресом и сделал запрос на освобождение памяти, не соответствующей началу блока памяти в аллокаторе. Чтобы отследить это, можно дополнительно хранить адреса занятых блоков, или же проверять адрес на кратность размеру блока в аллокаторе.

Применения пуллового аллокатора:

\begin{itemize}
	\item В игровых движках: для создания и уничтожения большого количества однотипных объектов (пули, враги, частицы) \cite{game++}.
	\item При обработке сетевых пакетов для многократного переиспользования буферов одинакового размера.
\end{itemize}

В своей книге Бьёрн Страуструп приводит пример класса Pool, который был разработан задолго до изобретения аллокаторов, однако принцип его работы был схож с описанным выше пулловым аллокатором. Данный класс формировал понятие пула элементов фиксированного размера, из которого пользователь мог быстро выделять память под объекты и освобождать её. Это был низкоуровневый тип, работающий с памятью напрямую и учитывающий необходимость в ручном выравнивании границ \cite{straustrup}.
